\chapter{The explicit formula as a deconvolution}\label{chap:es}\chaptermark{The explicit formula as deconvolution}

Chapters~\ref{chap:lo} and~\ref{chap:cx} write the Thorin measure at level $\theta$ as the Cauchy smoothing of the zeros at width $\theta-\tfrac12$ (under the Riemann hypothesis; in general the widths are $\theta-\sigma_\rho$), turn the passage from level $1$ to level $\theta$ into a deconvolution, and solve it by EM to an accuracy of about $10^{-3}$ in the positions of the zeros, with a certificate. This chapter looks at the bottom of that flow. At level $\tfrac12$, under the Riemann hypothesis, the Thorin measure is the counting measure of the zeros, the data are the level one density $u_1(v)=\pi^{-1}\Re(\xi'/\xi)(1+iv)$, and the deconvolution, written in Fourier variables, is Weil's explicit formula: the multiplier $e^{|\omega|/2}$ that undoes the Cauchy kernel is the factor $\sqrt n$ that turns the weights $\Lambda(n)/n$ of the level one spectrum into the weights $\Lambda(n)/\sqrt n$ of the formula. That identification says how ill posed the problem is. For a general measure the resolution $\delta$ costs, heuristically, the primes to $e^{\pi/\delta}$, or data accurate to $e^{-\pi/2\delta}$: exponential. For the measure at hand, which is a finite sum of atoms on every window, it costs much less, because the Fourier transform of a sum of Lorentzians is a sum of damped exponentials, and a sum of damped exponentials is recovered from its samples by Prony's method with an error that, for a fixed configuration of atoms, is proportional to the data error. The matrix pencil applied to the tapered Fourier transform of $u_1$ on a window of width $60$ returns, in the computation reported by the author, the ordinates and real parts of the central zeros to within $8\times10^{-10}$, from level one data alone, in seconds; a synthetic pair of zeros placed off the line at $0.85\pm60.089i$ is returned with real parts $0.85$ and $0.15$ to seven digits. So sparse methods solve the deconvolution many orders of magnitude better than EM on the same data, and the reason is structural: EM solves the problem on a dense grid, where it is exponentially ill posed, and the matrix pencil solves it in the model where it is polynomially ill posed. A posterior concentration theorem makes the recovery Bayesian, and the functional equation makes it singular: off line zeros come in pairs at $\tfrac12\pm\delta$, the data depend on $\delta$ only through $\delta^2$, and the real part of a pair is resolved at the rate $\tau^{1/2}$ when the data are accurate to $\tau$, with Bayes factors for the line against the pair that grow only like $\tau^{-1/2}$. None of this is a path to the hypothesis: the method reads real parts, it does not constrain them.

The chapter develops a source chapter on the explicit formula, sparse recovery and Bayesian detection of the zeros by Nick Polson \cite{P26es}, together with the computations done for this book. Its numbers were recomputed or checked here by scripts written from the mathematics alone (the scripts of the source were not used; the scripts of this chapter are \texttt{code/es\_*.py} and \texttt{code/ex\_zeros3e12.py}, with outputs \texttt{data/es\_*.json} and \texttt{data/ex\_zeros3e12.json}, listed in \texttt{code/README.md}). Every number printed below is one of those recomputations unless it is marked as the author's: a figure marked computed by the author's scripts was not rerun for this book, and where a recomputation differed from the source the recomputed value is printed. Several parts of the source are in the book already and are cited, not restated: the Lorentzian representation of $u_1$ and the Hadamard product (Proposition~\ref{prop:struct}), the criterion that the deconvolution has a positive solution (Theorem~\ref{tf:thm:deconv}), the explicit formula for even test functions read as the deconvolved functional (Theorem~\ref{dc:thm:weil}, Section~\ref{dc:sec:weil}), the EM, grid and certificate methods (Chapters~\ref{chap:lo} and~\ref{chap:cx}) and the survey of Cauchy deconvolution (Chapter~\ref{chap:cd}). What is new here is the form of Theorem~\ref{dc:thm:weil} for all compactly supported test functions, with its pole term, the taper theorem and its use with the matrix pencil, the computations, and the posterior theory with its singular case.

The notation is the book's, with the following warnings. A level $\theta$ is the real part of a point $s=\theta+iv$, and $u_\theta(v)=\pi^{-1}\Re(\xi'/\xi)(\theta+iv)$ is the density of Chapters~\ref{chap:lv} to~\ref{chap:cd}; in the notation of Chapter~\ref{chap:tf} (equation~\eqref{eq:rho}) it is $\rho_{\theta-1/2}(v)$, the ordinate being $v$ here. For a zero $\rho=\sigma_\rho+i\gamma_\rho$ the width is $a_\rho=1-\sigma_\rho\in(0,1)$ at level $1$. The Cauchy density is $P_a(x)=(a/\pi)/(a^2+x^2)$, $\mathfrak U_\theta=\sum_\rho\mathrm{Cauchy}(\gamma_\rho,\theta-\sigma_\rho)$ is the Thorin measure with density $\sum_\rho P_{\theta-\sigma_\rho}(v-\gamma_\rho)$ and no factor $\tfrac12$, as in Chapter~\ref{chap:cx}, so that $\mathfrak U_\theta=2\widetilde U_\theta$ with $\widetilde U_\theta$ of Theorem~\ref{lv:thm:semigroup}, $\gamma_E$ is Euler's constant and $\Lambda$ the von Mangoldt function. The source's Chapters 3 and 4 are Chapters~\ref{chap:lo} and~\ref{chap:cx} of the book. In the Bayesian sections $\tau$ is the standard deviation of each real coordinate of the noise. The word certified is used only for a computation in ball arithmetic whose sentence says what the certificate covers: the sign changes of Hardy's function $Z$ in Sections~\ref{es:sec:certpencil}, \ref{es:sec:above} and~\ref{es:sec:above13}. Everything else done in floating point is computed.

Status words follow the rest of the book: proved, computed, conditional, heuristic or open. The results are these.
\begin{enumerate}
\item The explicit formula of Guinand and Weil is the Fourier side of the deconvolution from level $1$ to level $\tfrac12$: for every test function in $C_c^\infty(\mathbb R)$, unconditionally, the zero side $\sum_\rho\int g\,e^{(\sigma_\rho-1/2)|\omega|}e^{-i\omega\gamma_\rho}d\omega$ equals the prime side with the pole term and the finite part of the Gamma factor; for even $g$ this is Theorem~\ref{dc:thm:weil}, and under the Riemann hypothesis it reads $\sum_\rho e^{-i\omega\gamma_\rho}=e^{|\omega|/2}(1+\widehat G)-\sum\Lambda(n)n^{-1/2}(\delta_{\log n}+\delta_{-\log n})$ (Theorem~\ref{es:thm:explicit}, proved; the reading is conditional). A Gaussian perturbation of the measure that data of accuracy $\varepsilon$ cannot see, of total variation $1$, exists as soon as $0.72\,e^{-\Omega/2}\le2\varepsilon$ (Proposition~\ref{es:prop:resolution}, proved); the consequence that resolution $\delta$ needs the primes to $e^{\pi/\delta}$ is a heuristic remark (Remark~\ref{es:rem:resolution}).
\item The tapered Fourier transform of $u_1$ on a window is a sum of damped exponentials, one per zero, with the pole $e^{\Delta(-a_\rho+i(\gamma_\rho-c))}$ carrying both coordinates of the zero, plus a remainder of order $e^{-\ell^2\omega^2/2}$ (Theorem~\ref{es:thm:taper}, proved, with no Gamma terms). The matrix pencil recovers the poles of an exact sum (Lemma~\ref{es:lem:pencil}, proved). Its error analysis is heuristic. At height $200$ the eleven zeros within $10$ of the centre are placed to within $8\times10^{-10}$ in both coordinates, and a synthetic pair off the line at $60.089$ is read to seven digits (Section~\ref{es:sec:computation}, computed by the author's scripts, not rerun).
\item Posterior concentration on a window holds at the rate $\tau$ for zeros with separated ordinates, under explicit hypotheses (Theorem~\ref{es:thm:concentration}, proved); the hypotheses do not hold at the settings of the computations (Remark~\ref{es:rem:settings}, computed), and concentration cannot prove the hypothesis, since a finite computation settles finitely many zeros.
\item The functional equation makes the model singular: with one amplitude common to the two members, a pair at $\tfrac12\pm\delta$ enters the data through $\delta^2$ (Lemma~\ref{es:lem:even}, proved). The posterior of $\delta^2$ concentrates at the rate $\tau$, so the real part is resolved at the rate $\tau^{1/2}$ when the line is the truth; no method, neural networks included, separates $\delta=0$ from $\delta=c\,(\tau/\|t_\perp\|)^{1/2}$, $t_\perp$ the projected tangent of the pair, with error probabilities summing to less than $1-c^2/2-O(\tau)$; and the matched filter and the likelihood ratio test attain the scale $\tau^{1/2}$ (Theorem~\ref{es:thm:singular}, proved for data simulated from the model in a window without aliasing, not for the runs on the zeros of $\zeta$).
\item Bayes factors with the right alternative, the pair and not a single zero: under the line the Bayes factor of the line against the pair is of order $\tau^{-1/2}$ in probability, under a fixed pair it is exponentially small, and against a single zero alternative it is of order $\tau^{-1}$ (Theorem~\ref{es:thm:bf}, proved in the common amplitude model). For the zero at $201.264752$, at an assumed noise level, the computed $\log_{10}$ of the Bayes factor is $6.1$ against the pair and $11.8$ against the single zero alternative: the functional equation takes roughly the square root of the evidence (computed on an assumed design, that of the detection experiment, an illustration).
\item An adaptive prior, the sine kernel process on the distinct ordinates at the Riemann and von Mangoldt density and a spike and slab on the pairs for the real parts, is a heuristic choice resting on a conjecture about the ordinates; its properties used here are proved (Remarks~\ref{es:rem:singular} and~\ref{es:rem:dpp}). A detection experiment compares the pencil and a plug in matched filter with the benchmark of Theorem~\ref{es:thm:singular}(iii) (computed).
\item The section on the open problems of the source works them: the certification of the pencil's output by sign changes and a winding number (logic proved, eleven zeros near height $200$ computed), close pairs, the levels $\tfrac34$ and $0.6$ and the signed level $\theta$ transform (Proposition~\ref{es:prop:level}, proved), the level as a cost parameter, and many windows (173 zeros on $[100,400]$ from 31 windows, computed by the author).
\item A window at height $c=3\times10^{12}+100$ lies inside the range where Platt and Trudgian \cite{PT21} verify the hypothesis, $0<\gamma\le3\,000\,175\,332\,800$, so the run is a demonstration at a height near the limit of the verification: the data are values of $\xi'/\xi$ at $\Re s=0.6$ computed in ball arithmetic, no primes are used and no zero unknown before is found. The pencil returns $52$ components; an independent scan of Hardy's function finds exactly $52$ certified sign changes there, and the one by one comparison with the pencil's ordinates has not been made. The selection of components in the source's runs presumes the real part $\tfrac12$ (Section~\ref{es:sec:above}, computed).
\item A window at $c=10^{13}+100$, above the verified height by a factor $3.3$, is run with the book's own pencil, whose selection uses the width of each pole: $53$ components, real parts within $5\times10^{-9}$ of $\tfrac12$ (computed, an uncertified estimate), every ordinate within $10^{-7}$ of one of $53$ certified sign changes of $Z$ and no certified sign change without a pencil ordinate. The count of zeros in the window and the real parts are not certified, and nothing is proved about any zero outside the $53$ (Section~\ref{es:sec:above13}).
\end{enumerate}
The open problems close the chapter.
